\documentclass[12pt,a4paper]{article} % ใช้ XeLaTex compile
\usepackage{DSIpreamble}
\usepackage{hyperref} % เพื่อให้ใช้ url ได้
\usepackage{amsmath} % เพื่อให้เขียนเวกเตอร์และเมทริกซ์ได้

% แก้ชื่อและหัวข้อตรงนี้
\author{เจน โดว์ 6103456789 \\ John Doe 6203456789}
\title{วิธีการใช้ \LaTeX ในการเขียน Assignment}

\begin{document}
\maketitle

% ใส่เนื้อหาตรงนี้
แนะนำวิธีการใช้ \LaTeX สำหรับทำงาน Assignment ในวิชา วข.205	พีชคณิตเชิงเส้นพื้นฐานสำหรับการวิเคราะห์ข้อมูล		
DSI205	Basics of Linear Algebra for Data Analytics

\section{ติดตั้ง Latex บน PC}
\begin{enumerate}
\item ดาวน์โหลด และ Install \code{MiKTex} จาก \url{https://miktex.org/} ซึ่งมาพร้อมกับ \code{TeXworks} ที่เป็น Editor \\
\code{MiKTex} เป็น Compiler ที่จะแปลงโค้ด Latex เป็น PDF ส่วน \code{TeXworks} เป็นแค่ Editor ที่ทำให้เขียนโค้ด Latex ได้สะดวก
\item ดาวน์โหลดโค้ดตัวอย่างจาก \url{https://www.overleaf.com/read/hyhynwsnpjdx}  \\
ให้ไปที่ Menu > Download > Source
\item เปิดไฟล์ \code{.tex} ด้วย \code{TeXworks} แล้วเปลี่ยน Compiler เป็น \textbf{XeLaTex} ที่ Menu
\item แล้วก็ปุ่มลูกศรสีเขียว เพื่อ Compile/Typeset
\item กด Install หากโปรแกรมแจ้งว่าต้องการติดตั้ง Library
\item เสร็จแล้วก็จะได้ PDF ไฟล์
\end{enumerate}


\section{ใช้ Overleaf เขียน Online}
คุณสามารถเขียน Latex ออนไลน์ได้ บนเว็บ \emph{Overleaf} \url{https://www.overleaf.com} \\
คำเตือน: เว็บ \emph{Overleaf} จะ Compile ค่อนข้างช้า
\begin{enumerate}
\item สมัครสมาชิกที่ \href{https://www.overleaf.com}{Overleaf}
\item ไปที่ \url{https://www.overleaf.com/latex/templates/dsi205-template/tsnyrzxnngdw} และ Open as Template (ต้องล็อกอินเว็บ Overleaf ก่อน)
\item ไปที่ Menu > Compiler เปลี่ยน Compiler เป็น \textbf{XeLaTex}
\item กด Recompile หรือตั้งให้เป็น Auto-Compile (โดยกดลูกศรข้าง ๆ แล้วเลือก Auto Compile เป็น On)
\item เมื่อแก้ไฟล์เสร็จแล้ว ให้ Recompilie แล้วดาวน์โหลด PDF โดยเลือก Menu > Download > PDF
\end{enumerate}

\section{การเขียนสมการ}
ต้องใช้ Math Mode ซึ่งมี 3 วิธีด้วยกัน

\subsection{Inline mode}
สมการจะถูกแสดงอยู่ในบรรทัดเดียวกับข้อความ เราสามารถใช้เครื่องหมาย $\backslash( \dots \backslash)$, $\$ \dots \$$, หรือ $\backslash$begin\{math\} $\dots \backslash$end\{math\} เช่น

\code{\$E = mc\^2\$} เป็น $E = mc^2$

\code{$\backslash($ E = mc\^2 $\backslash)$} เป็น \( E = mc^2 \)

\subsection{Displayed Mode แบบ Unnumbered}
สมการจะขึ้นบรรทัดใหม่ แต่ไม่มีตัวเลขกำกับ เราสามารถใช้เครื่องหมาย $\backslash[ \dots \backslash]$ หรือ \$\$ $\dots$ \$\$ เช่น

\code{$\backslash[$ $\backslash$bar\{x\} = $\backslash$sum x\_i $\backslash]$} เป็น
\[ \bar{x} = \sum x_i \]

\code{\$\$ $\backslash$bar\{x\} = $\backslash$sum x\_i \$\$} เป็น
$$ \bar{x} = \sum x_i $$

\subsection{Displayed Mode แบบ Numbered}
สมการจะขึ้นบรรทัดใหม่ และมีตัวเลขกำกับซึ่งจะเพิ่มขึ้นเรื่อย ๆ ตามจำนวนสมการในเอกสาร เราสามารถใช้เครื่องหมาย $\backslash$begin\{equation\} $\dots \backslash$end\{equation\}

เช่น \\
$\backslash$begin\{equation\} \\
I = $\backslash$frac\{V\}\{R\} \\
$\backslash$end\{equation\}

เป็น
\begin{equation}
I = \frac{V}{R}
\end{equation}

\section{การเขียนสมการหลายบรรทัด}
ใช้ \code{aligned} ใน \code{equation}
\begin{equation}
\begin{aligned}
a & = 2 \times 2 \times b \times c \\
a & = 4 \times b \times c \\
a & = 4bc
\end{aligned}
\end{equation}

\section{การเขียนเวกเตอร์ และเมทริกซ์}
ใช้แพกเกจ \code{amsmath} และคำสั่ง \code{bmatrix} ใน \code{equation}
\begin{equation}
{\mathbf v} = \begin{bmatrix}1 \\ 5 \\ -4 \\ 0\end{bmatrix}
\end{equation}

\begin{equation}
A = \begin{pmatrix}1 & 2 & -3 \\ 4 & 0 & 1\end{pmatrix}
\end{equation}

\begin{equation}
B = \begin{bmatrix}1 & 2 & -3 \\ 4 & 0 & 1\end{bmatrix}
\end{equation}

\begin{equation}
c = \det A = \begin{vmatrix}-2 & 0 & 1 \\ 5 & -4 & 2 \\ 3 & 1 & 8\end{vmatrix}
\end{equation}

\end{document}